\chapter{Synchronization}\label{ch:sync}

A receiver does not know when symbols start, where the frame begins, or the exact carrier
frequency and phase. \Cref{fig:sync-chain} shows the synchronized single-carrier receiver of
\file{phy/sync.py}, which recovers these quantities in three stages.

\begin{figure}[htbp]
\centering
\begin{tikzpicture}[node distance=4mm and 5mm]
  \node[blockbox] (map) {mapper};
  \node[blockbox,right=of map] (pre) {preamble};
  \node[blockbox,right=of pre] (ps) {pulse\\shaper};
  \node[warmbox,right=of ps,minimum width=26mm] (ch) {delay $\tau$, CFO $\Delta f$,\\phase noise, AWGN};
  \node[blockbox,right=of ch] (mf) {matched\\filter};
  \node[greenbox,below=9mm of mf] (tr) {timing\\recovery};
  \node[greenbox,left=of tr] (fs) {frame\\sync};
  \node[greenbox,left=of fs] (cr) {carrier\\recovery};
  \node[blockbox,left=of cr] (dm) {demapper};
  \draw[sig] (map)--(pre); \draw[sig] (pre)--(ps); \draw[sig] (ps)--(ch); \draw[sig] (ch)--(mf);
  \draw[sig] (mf)--node[right,lbl]{$N_{\mathrm{sps}}\ge2$}(tr); \draw[sig] (tr)--node[above,lbl]{1 sps}(fs);
  \draw[sig] (fs)--(cr); \draw[sig] (cr)--(dm);
  \node[note,below=1mm of tr] {fractional $\tau$};
  \node[note,below=1mm of fs] {integer offset,\\coarse $\Delta f$, phase};
  \node[note,below=1mm of cr] {residual phase};
\end{tikzpicture}
\caption{The synchronized receiver of the \emph{Synchronized QPSK receiver} scenario.}
\label{fig:sync-chain}
\end{figure}

\begin{figure}[htbp]
\centering
\begin{tikzpicture}
\begin{groupplot}[group style={group size=4 by 1,horizontal sep=4mm,y descriptions at=edge left},
  constaxis,width=4.1cm,height=4.1cm]
\nextgroupplot[title={nominal sampling}]
  \addplot[scatterpts=black!50] table[x=i,y=q]{\datafile{sync_raw}};
\nextgroupplot[title={after timing rec.}]
  \addplot[scatterpts=accent] table[x=i,y=q]{\datafile{sync_tr}};
\nextgroupplot[title={after frame sync}]
  \addplot[scatterpts=accent4] table[x=i,y=q]{\datafile{sync_fs}};
\nextgroupplot[title={after the PLL}]
  \addplot[scatterpts=accent3] table[x=i,y=q]{\datafile{sync_cr}};
\end{groupplot}
\end{tikzpicture}
\caption{QPSK through a delay of 9.37 samples, a 3\,kHz CFO at 1\,MBd, 10\,Hz phase noise and
12\,dB $E_b/N_0$. Panels: sampled at the nominal instants (ISI and rotation), after Gardner
timing recovery (rotating ring), after preamble-based CFO and phase correction (residual
spread), after the decision-directed PLL.}
\label{fig:sync-const}
\end{figure}

% ------------------------------------------------------------------
\section{The preamble}
\blk{preamble} prepends $L$ known QPSK symbols made of two identical halves of length
$h=L/2$, the Schmidl--Cox structure~\cite{schmidl1997}, and appends \prm{guard} zero symbols.
The repetition makes the timing metric insensitive to the CFO and yields a frequency estimate
(\cref{sec:frame-sync}). The guard gives the receiver slack for unknown delays.

% ------------------------------------------------------------------
\section{Symbol timing recovery}
After the matched filter the waveform is sampled at $N_{\mathrm{sps}}$ samples per symbol, but
the optimum instants $mT+\tau$ fall between samples. Timing recovery has three parts:
\begin{itemize}
  \item an \emph{interpolator} that evaluates the waveform at arbitrary instants;
  \item a \emph{timing error detector} (TED);
  \item a \emph{loop filter} that turns detector outputs into corrections.
\end{itemize}

\subsection{Interpolation}
A cubic Lagrange (Farrow) interpolator evaluates $x(i+\mu)$ for $0\le\mu<1$ from four
neighbours:
\begin{equation}
  x(i+\mu)=c_0+c_1\mu+c_2\mu^2+c_3\mu^3,\quad
  \begin{aligned}
  c_0&=x_0, & c_1&=-\tfrac13x_{-1}-\tfrac12x_0+x_1-\tfrac16x_2,\\
  c_2&=\tfrac12x_{-1}-x_0+\tfrac12x_1, & c_3&=-\tfrac16x_{-1}+\tfrac12x_0-\tfrac12x_1+\tfrac16x_2 .
  \end{aligned}
\end{equation}
\codeexcerpt{cubic}{Farrow-structure cubic interpolation (\code{cubic_interpolate}).}

\subsection{Feedforward start: the Oerder--Meyr estimator}
A linearly modulated signal is cyclostationary. The mean of $|x(t)|^2$ is periodic in $T$, and
its first Fourier coefficient has a phase proportional to the timing offset. Over $K$ symbols,
\begin{equation}
  \hat\varepsilon=-\frac{1}{2\pi}\arg\sum_{k=0}^{KN_{\mathrm{sps}}-1}|x[k]|^2e^{-j2\pi k/N_{\mathrm{sps}}}
  \qquad(\text{in symbols, modulo }1),
\end{equation}
which needs $N_{\mathrm{sps}}\ge3$~\cite{oerder1988}. The block uses it to initialize the
loop, which removes the long acquisition (and the risk of hang-up) of a loop started from an
arbitrary phase.
\codeexcerpt{om_estimate}{Oerder--Meyr initial estimate.}

\subsection{The Gardner detector}
Gardner's TED~\cite{gardner1986} uses two samples per symbol, the current strobe and the
midpoint between strobes:
\begin{equation}
  e[m]=\Re\Big\{x\conj\big((m-\tfrac12)T+\hat\tau\big)\,\big[x(mT+\hat\tau)-x((m-1)T+\hat\tau)\big]\Big\}.
\end{equation}
It is independent of the carrier phase and works before carrier recovery. Its mean response
to a timing error, the \emph{S-curve}, is zero at the correct instant and has slope $K_d$
there (\cref{fig:scurve}). With the convention of the code, $e>0$ means the strobes are late.

\begin{figure}[htbp]
\centering
\begin{tikzpicture}
\begin{axis}[width=9.5cm,height=5.4cm,xlabel={timing error $\tau/T$},ylabel={$\E\,e$},
  xmin=-0.5,xmax=0.5,legend pos=north west]
  \addplot[th=accent] table[x=tau,y=b20]{\datafile{gardner_scurve}};
  \addplot[th=accent3] table[x=tau,y=b35]{\datafile{gardner_scurve}};
  \addplot[th=accent2] table[x=tau,y=b60]{\datafile{gardner_scurve}};
  \legend{$\beta=0.2$,$\beta=0.35$,$\beta=0.6$}
\end{axis}
\end{tikzpicture}
\caption{Measured S-curves of the Gardner detector for QPSK with RRC pulses (8 samples per symbol,
matched-filtered output of the real blocks, interpolated with the same cubic interpolator). The
slope at the origin grows with the roll-off; for $\beta=0.35$ it is close to 1, the default
\prm{detector_gain}.}
\label{fig:scurve}
\end{figure}

\subsection{The second-order loop}\label{sec:loop}
\begin{figure}[htbp]
\centering
\begin{tikzpicture}[node distance=5mm and 9mm]
  \node[blockbox] (ted) {detector\\$K_d$};
  \node[blockbox,right=of ted] (kp) {$K_p$};
  \node[blockbox,below=of kp] (ki) {$K_i\dfrac{1}{1-z^{-1}}$};
  \node[sumnode,right=of kp] (s) {$+$};
  \node[blockbox,right=of s] (nco) {NCO\\$\dfrac{1}{1-z^{-1}}$};
  \node[blockbox,right=of nco] (int) {interpolator};
  \draw[sig] (ted) -- node[above,lbl]{$e$} (kp);
  \draw[sig] (ted.south) |- (ki);
  \draw[sig] (kp) -- (s); \draw[sig] (ki) -| (s);
  \draw[sig] (s) -- (nco); \draw[sig] (nco) -- node[above,lbl]{$\hat\tau$} (int);
  \draw[sig] (int.north) -- ++(0,6mm) -| node[pos=0.25,above,lbl]{strobes} (ted.north);
\end{tikzpicture}
\caption{Proportional-plus-integral loop used for timing (and, with a phase detector, for the
PLL).}
\end{figure}
A proportional-plus-integral loop filter makes the loop second order. It tracks a constant
offset with zero error and a frequency (clock drift) with zero steady-state error. The gains
follow from the normalized noise bandwidth $B_nT$ and the damping $\zeta$~\cite{mengali1997}:
\begin{equation}
  \theta=\frac{B_nT}{\zeta+\frac{1}{4\zeta}},\qquad
  K_p=\frac{4\zeta\theta}{(1+2\zeta\theta+\theta^2)K_d},\qquad
  K_i=\frac{4\theta^2}{(1+2\zeta\theta+\theta^2)K_d} .
\end{equation}
\codeexcerpt{loop_gains}{Second-order loop gains.}
\codeexcerpt{gardner_loop}{The timing loop: interpolate, detect, filter, advance.}
A small $B_nT$ averages noise better but tracks more slowly; $0.01$ (the default) is a common
compromise.

% ------------------------------------------------------------------
\section{Frame synchronization and coarse CFO}\label{sec:frame-sync}
At symbol rate, with a residual CFO of $\epsilon$ cycles per symbol, the two halves of the
preamble differ only by the phase $2\pi\epsilon h$. The Schmidl--Cox metric is evaluated at
every candidate start $d$:
\begin{equation}
  P(d)=\sum_{k=0}^{h-1}r\conj[d+k]\,r[d+k+h],\qquad
  M(d)=\frac{|P(d)|^2}{\big(\sum_k|r[d+k+h]|^2\big)^2},\qquad
  \hat\epsilon=\frac{\arg P(\hat d)}{2\pi h}.
\end{equation}
The estimate is unambiguous for $|\epsilon|<1/(2h)$. Because $M(d)$ has a plateau, its maximum is
used only for the CFO estimate. After derotation, the frame start is located by
cross-correlation with the known preamble,
$\hat d=\argmax_d\big|\sum_kr[d+k]\,p\conj[k]\big|$. The CFO is then refined at the exact
position, and the common phase is taken from the correlation peak. Finally the payload symbols
are extracted.
\codeexcerpt{frame_sync}{Schmidl--Cox CFO, correlation timing and refinement (\blk{frame_sync}).}

% ------------------------------------------------------------------
\section{Carrier phase tracking}
The residual phase (CFO residue and phase noise) is tracked by a decision-directed PLL. For
$y=s\,e^{j\Delta\phi}+w$ and a correct decision $\hat s$,
\begin{equation}
  e=\frac{\Im\{y\,\hat s\conj\}}{|\hat s|^2}=\sin\Delta\phi+\text{noise}\approx\Delta\phi ,
\end{equation}
so the detector gain is 1. The same second-order loop as above updates the phase estimate
every symbol. Like every decision-directed scheme it has an $M$-fold phase ambiguity, which the
preamble phase estimate resolves beforehand.
\codeexcerpt{dd_pll}{Decision-directed PLL (\blk{carrier_recovery}).}

\begin{pitfall}
The PLL bandwidth must exceed the phase-noise linewidth relative to the symbol rate. It must
also be small enough that decision errors at low SNR do not make it slip cycles. A cycle slip
rotates all subsequent decisions by $2\pi/M$ and appears as a burst of errors until the next
preamble. Differential encoding or pilot symbols are the classic remedies.
\end{pitfall}
